\chapter{Rapidly-Exploring Random Trees}
\label{ch:ch16}
% Function names for pseudocode are prefixed with "Rrt" to stay unique
% across the book.
\SetKwFunction{RrtGrow}{RRT}
\SetKwFunction{RrtSampleBiased}{Sample}
\SetKwFunction{RrtNearestVertex}{Nearest}
\SetKwFunction{RrtSteerEta}{Steer}
\SetKwFunction{RrtSegmentFree}{CollisionFree}
\SetKwFunction{RrtExtendTree}{Extend}
\SetKwFunction{RrtConnectGreedy}{Connect}
\SetKwFunction{RrtConnectPlan}{RRT-Connect}
\SetKwFunction{RrtExtractPath}{ExtractPath}
\SetKwFunction{RrtKinoGrow}{KinodynamicRRT}
\SetKwFunction{RrtIntegrate}{Integrate}
\SetKwFunction{RrtShortcutPath}{Shortcut}
\SetKwFunction{RrtPrunePath}{Prune}
\SetKwFunction{RrtSwap}{Swap}

\begin{objectives}
\begin{itemize}
  \item Explain why grids fail in large or high-dimensional continuous spaces, and what a sampling-based planner does instead.
  \item State the four primitives of \rrt exactly (\textsc{Sample} with goal bias, \textsc{Nearest}, \textsc{Steer} with step size $\eta$, \textsc{CollisionFree} by subdivision) and trace \rrt by hand on a small map.
  \item Explain the Voronoi bias, prove that \rrt is probabilistically complete, and explain why it is not optimal.
  \item Use \rrt-Connect, kinodynamic \rrt and path shortcutting, and estimate the cost of nearest-neighbor search and collision checking.
  \item Implement \rrt in \python for 2D and 3D box worlds, run it in a 3D obstacle field, and decide when to use grid search and when to sample.
\end{itemize}
\end{objectives}

% ---------------------------------------------------------------------
\section{Why this matters for a drone swarm}
\label{sec:ch16-motivation}

Every planner of \cref{ch:ch03,ch:ch04,ch:ch05,ch:ch06} and every multi-agent
planner of \cref{ch:ch07,ch:ch08,ch:ch09,ch:ch10,ch:ch11} searches a graph
that somebody built first, usually a grid. On a $100\times100$ city map with
one-meter cells that is a fine choice. Now let the drone fly through a
$100\times100\times50$~m volume between buildings, cranes and trees, with
a required accuracy of ten centimeters. The grid has $5\times10^{8}$ cells
(\cref{tab:ch16-gridsizes}); it needs half a gigabyte to store one byte per
cell, and \astar has to touch a good fraction of it before it reaches the
goal. Add the velocity of the drone to the state, because a fast drone
cannot turn a corner in place, and the grid has $5\times10^{14}$ cells: no
computer
holds it. This is the \textbf{curse of
dimensionality}\index{curse of dimensionality}: the number of cells grows as
$N^{d}$ with the resolution $N$ per axis and the dimension $d$, and so does
the work of any method that must visit them one by one.

\begin{table}[tb]
\centering\small
\caption{The curse of dimensionality. Sizes of a uniform grid over a
$100\times100$~m map, a $100\times100\times50$~m volume, and the same volume
with the velocity ($\pm5$~m/s per axis) added to the state. The last column
is the number of neighbors of a cell in the fully connected lattice,
$3^{d}-1$. A tree of $10^{4}$ vertices with three coordinates each takes
about $240$~kB.}
\label{tab:ch16-gridsizes}
\begin{tabular}{@{}lllrrr@{}}
\toprule
Space & Resolution & Cells per axis & Cells & Memory & Neighbors\\
\midrule
2D map & 0.5~m & 200 & $4.0\times10^{4}$ & 40~kB & 8\\
2D map & 0.1~m & 1000 & $1.0\times10^{6}$ & 1~MB & 8\\
3D volume & 0.5~m & 200, 200, 100 & $4.0\times10^{6}$ & 4~MB & 26\\
3D volume & 0.1~m & 1000, 1000, 500 & $5.0\times10^{8}$ & 500~MB & 26\\
6D position+velocity & 0.5~m, 0.5~m/s & as above, 20 & $3.2\times10^{10}$ & 32~GB & 728\\
6D position+velocity & 0.1~m, 0.1~m/s & as above, 100 & $5.0\times10^{14}$ & 500~TB & 728\\
\bottomrule
\end{tabular}
\end{table}

Sampling-based planning\index{sampling-based planning} turns the problem
around. Instead of enumerating the space it draws random points from it,
keeps the collision-free ones and connects them by short collision-free
segments. The memory is proportional to the number of samples, not to the
size of the space, and the same code runs in two, three or twelve
dimensions. The rapidly-exploring random tree (\rrt) of
\textcite{lavalle1998rrt} is the simplest and most used member of this
family: it grows a tree from the start, one random sample at a time, and
stops when the tree reaches the goal. In the hybrid architecture of
\cref{ch:ch24} it is the continuous planner that takes over when the global
grid is too coarse for the local geometry or too large to build
(\cref{sec:ch16-drone}). Week~8 of the study plan (\cref{ch:appA}) asks you
to implement a sampling-based planner in a 3D obstacle field and to compare
it with grid \astar; this chapter and \cref{ch:ch17} are its material.

\Cref{sec:ch16-intuition} gives the idea, \cref{sec:ch16-problem} the exact
primitives, \cref{sec:ch16-algorithm} the algorithm and
\cref{sec:ch16-example} a worked example. \Cref{sec:ch16-properties} proves
probabilistic completeness and shows why \rrt is not optimal.
\Cref{sec:ch16-variants} covers \rrt-Connect, kinodynamic \rrt, shortcutting
and planning in 3D, \cref{sec:ch16-implementation} the collision checks,
nearest-neighbor search and parameters, and \cref{sec:ch16-choice} the
choice between grid search and sampling.

% ---------------------------------------------------------------------
\section{The idea in plain words}
\label{sec:ch16-intuition}

Imagine throwing darts at a map. The tree, which starts as the single start
point, reacts to each dart in the same way: it finds the vertex closest to
the dart, moves from that vertex a fixed distance $\eta$ toward the dart,
and, if that short step does not cross an obstacle, adds the end of the step
as a new vertex with the old vertex as its parent. The dart itself is thrown
away. After a few hundred darts the tree has branches everywhere, and as soon
as one branch enters a small region around the goal, the chain of parents
from that vertex back to the root is a collision-free path.

Why does this explore so quickly, without any heuristic pointing to the
goal? Because of the way the nearest vertex is chosen.
\Cref{fig:ch16-voronoi} shows a small tree in an empty square together with
the \textbf{Voronoi region}\index{Voronoi region} of every vertex, the set of
points closer to this vertex than to any other. A dart that lands in the
region of vertex $\vect{v}$ makes $\vect{v}$ the parent of the new vertex, so with uniform
darts the probability that $\vect{v}$ is \emph{selected} equals the area of its
region divided by the area of the map; whether the step is then accepted is
up to the collision check. The tips of branches that point into empty space
own the largest regions, so they are selected most often, and every accepted
extension pushes the tree further into the empty space. In the figure the
largest region covers $20.6\,\%$ of the square and the three largest
together $59.3\,\%$, while an even share would be $6.7\,\%$ per vertex; the
smallest region, deep inside the tree, covers $0.7\,\%$. This is the
\textbf{Voronoi bias}\index{Voronoi bias}\index{RRT!Voronoi bias} of \rrt:
the tree is pulled toward the largest unexplored regions until they have
been filled.

\begin{figure}[tb]
  \centering
  \inputfigure{ch16/voronoi}
  \caption{The Voronoi bias. A 15-vertex \rrt in an empty $10\times10$ square
  (step size $\eta=1$) and the Voronoi regions of its vertices. The vertex
  $\state_{\mathrm{near}}$ at the tip of the longest branch owns the orange region,
  $20.6\,\%$ of the square, so one uniform sample in five extends it; the
  sample $\state_{\mathrm{rand}}$ shown does, and the orange step of length $\eta$
  pushes the branch further into the empty part of the map. The small regions
  inside the tree are almost never chosen.}
  \label{fig:ch16-voronoi}
\end{figure}

\begin{keyidea}
Throw random darts at the space; each dart pulls the nearest vertex of the
tree one step of length $\eta$ toward it. A vertex is selected with
probability equal to the size of its Voronoi region, and pulled if the step
is collision-free, so the tree rushes into
the largest unexplored regions without any heuristic. Nothing, however,
makes its paths short: \rrt finds \emph{a} path, not the best one.\index{RRT!key idea}
\end{keyidea}

% ---------------------------------------------------------------------
\section{Problem statement and notation}
\label{sec:ch16-problem}

\subsection{The space and the query}

\Cref{ch:ch02} introduced the configuration space $\Cspace$, its obstacle
region $\Cobs$ and its free space $\Cfree$. The sampling-based planning
literature writes the same objects as $\mathcal{X}$, $\Xobs$ and $\Xfree$,
because the space is often a \emph{state} space that contains velocities as
well as positions (\cref{sec:ch16-kinodynamic}); this chapter and
\cref{ch:ch17} follow that usage. For a drone planned as a sphere of radius
$r$, a configuration is the position of its center, $\state\in\R^{d}$ with $d=2$
or $3$, and $\Xobs$ is the union of the obstacles inflated by $r$ (the
Minkowski sum of \cref{ch:ch02}).

\begin{definition}[Sampling-based planning problem]
\label{def:ch16-problem}
Let $\mathcal{X}\subset\R^{d}$ be a bounded box, the \textbf{configuration
space}\index{configuration space!in sampling-based planning}, let
$\Xobs\subset\mathcal{X}$ be a closed set, the \textbf{obstacle region}, and
let $\Xfree=\mathcal{X}\setminus\Xobs$ be the \textbf{free
space}\index{free space}. A query consists of an initial configuration
$\state_{\mathrm{init}}\in\Xfree$ and a \textbf{goal region}\index{goal region}
$X_{\mathrm{goal}}\subseteq\Xfree$; in this book the goal region is the ball
$X_{\mathrm{goal}}=\set{\state : \norm{\state-\state_{\mathrm{goal}}}\le r_{\mathrm{goal}}}$
of radius $r_{\mathrm{goal}}>0$ around a goal point $\state_{\mathrm{goal}}$.
A \textbf{solution} is a continuous path
$\sigma\colon[0,1]\to\Xfree$ with $\sigma(0)=\state_{\mathrm{init}}$ and
$\sigma(1)\in X_{\mathrm{goal}}$. The planners of this chapter return
polygonal paths, sequences of configurations $\state_0=\state_{\mathrm{init}},\state_1,\dots,\state_k$
whose connecting segments lie in $\Xfree$; the \textbf{length} of such a path is
$\sum_{i}\norm{\state_{i}-\state_{i-1}}$, the path length of \cref{ch:ch02}.
\end{definition}

The tree $\mathcal{T}=(V,E)$ built by \rrt has a vertex set $V\subset\Xfree$ rooted at
$\state_{\mathrm{init}}$ and one directed edge $(\state_{\mathrm{parent}},\state)$ per
vertex except the root, each a collision-free segment; the path from the
root to any vertex is therefore a solution as soon as the vertex lies in
$X_{\mathrm{goal}}$.

\subsection{The four primitives}
\label{sec:ch16-primitives}

\rrt is defined by four operations, and every variant in this chapter and the
next changes one of them. \Cref{fig:ch16-primitives} shows them together.

\begin{definition}[Sample, Nearest, Steer, CollisionFree]
\label{def:ch16-primitives}
Let $\eta>0$ be the \textbf{step size}\index{step size}\index{RRT!step size} and
$p_{\mathrm{goal}}\in[0,1)$ the \textbf{goal bias}\index{goal bias}\index{RRT!goal bias}.
\begin{itemize}
  \item $\textsc{Sample}()$\index{RRT!Sample} returns $\state_{\mathrm{goal}}$ with probability
        $p_{\mathrm{goal}}$ and otherwise a configuration drawn uniformly at random
        from $\mathcal{X}$ (not from $\Xfree$: the sample may lie inside an obstacle,
        the collision check below takes care of it). Successive samples are
        independent. (\cref{sec:appB-probability} reviews random vectors,
        densities and expectations, the probability the completeness argument
        of \cref{sec:ch16-properties} rests on.)
  \item $\Nearest(V,\state)=\argmin_{\vect{v}\in V}\norm{\vect{v}-\state}$\index{RRT!Nearest} is the vertex of the
        tree closest to $\state$ in the Euclidean norm, ties broken arbitrarily.
  \item $\Steer(\state_{\mathrm{near}},\state_{\mathrm{rand}})$\index{RRT!Steer} moves from
        $\state_{\mathrm{near}}$ toward $\state_{\mathrm{rand}}$ by at most $\eta$:
        \begin{equation}
        \Steer(\state_{\mathrm{near}},\state_{\mathrm{rand}})=
        \begin{cases}
          \state_{\mathrm{rand}} & \text{if } \norm{\state_{\mathrm{rand}}-\state_{\mathrm{near}}}\le\eta,\\[2pt]
          \state_{\mathrm{near}}+\eta\,\dfrac{\state_{\mathrm{rand}}-\state_{\mathrm{near}}}{\norm{\state_{\mathrm{rand}}-\state_{\mathrm{near}}}} & \text{otherwise.}
        \end{cases}
        \label{eq:ch16-steer}
        \end{equation}
  \item $\CollisionFree(\vect{a},\vect{b})$\index{RRT!CollisionFree} is true if and only if the closed
        segment $[\vect{a},\vect{b}]=\set{\vect{a}+t(\vect{b}-\vect{a}): t\in[0,1]}$ lies in $\Xfree$.
\end{itemize}
\end{definition}

\begin{figure}[tb]
  \centering
  \inputfigure{ch16/primitives}
  \caption{The primitives of one iteration. The sample $\state_{\mathrm{rand}}$
  selects the nearest vertex $\state_{\mathrm{near}}$; \textsc{Steer} moves at most
  $\eta$ toward the sample and proposes $\state_{\mathrm{new}}$;
  \textsc{CollisionFree} tests points spaced $\delta$ apart along the new
  segment against the obstacle inflated by $r+\delta/2$ (dashed), which is
  conservative by \cref{thm:ch16-resolution}. The second attempt, from
  $\state_{\mathrm{new}}$ toward $\state'_{\mathrm{rand}}$, is rejected because its
  second test point falls inside the inflated region; no vertex is added.}
  \label{fig:ch16-primitives}
\end{figure}

A segment has infinitely many points, so \textsc{CollisionFree} is
implemented by testing finitely many of them. The following proposition says
how to do this without ever missing an obstacle. Write $\dist(\state,\mathcal{O})$
for the Euclidean distance from $\state$ to the workspace obstacles $\mathcal{O}$;
for axis-aligned boxes and spheres it has a closed form
(\cref{sec:ch16-collision}).

\begin{proposition}[Collision checking by subdivision]
\label{thm:ch16-resolution}
Let $r\ge0$ be the robot radius, so that a point $\state$ is free if and only if
$\dist(\state,\mathcal{O})>r$, and let $\delta>0$ be the \textbf{check
resolution}\index{collision checking!resolution}\index{collision checking!by subdivision}.
Let $\vect{p}_0=\vect{a},\vect{p}_1,\dots,\vect{p}_n=\vect{b}$ be points on the segment $[\vect{a},\vect{b}]$ with
$\norm{\vect{p}_{i}-\vect{p}_{i-1}}\le\delta$ for all $i$. If $\dist(\vect{p}_i,\mathcal{O})>r+\delta/2$
for every $i$, then every point $\vect{p}$ of $[\vect{a},\vect{b}]$ has $\dist(\vect{p},\mathcal{O})>r$, that is,
the segment is collision-free.
\end{proposition}
\begin{proof}
Take any $\vect{p}\in[\vect{a},\vect{b}]$. It lies between two consecutive test points, so its
distance to the nearer of them, $\vect{p}_i$ say, is at most $\delta/2$. The function
$\state\mapsto\dist(\state,\mathcal{O})$ is $1$-Lipschitz (an infimum of distances,
each of which changes by at most $\norm{\vect{p}-\vect{p}_i}$ when $\vect{p}_i$ is replaced by
$\vect{p}$), so
$\dist(\vect{p},\mathcal{O})\ge\dist(\vect{p}_i,\mathcal{O})-\norm{\vect{p}-\vect{p}_i}>r+\delta/2-\delta/2=r$.
\qedhere
\end{proof}

The test rejects some free segments, those that pass within $r+\delta/2$ of
an obstacle, and needs $\lceil\norm{\vect{b}-\vect{a}}/\delta\rceil+1$ point tests. It
never accepts a colliding segment, whatever the shape or thickness of the
obstacles, which is what makes it safe for a drone. Without the inflation by
$\delta/2$ a wall thinner than $\delta$ can slip between two test points
(\cref{sec:ch16-pitfalls}); the self-test of \code{ch16\_rrt.py} checks
exactly this case, a wall of thickness $0.02$ with $\delta=0.05$, at forty
segment angles through a point inside the wall and at forty parallel
crossings of it.

% ---------------------------------------------------------------------
\section{The algorithm}
\label{sec:ch16-algorithm}

\begin{algorithm}[htb]
\caption[\rrt with goal bias]{Rapidly-exploring random tree with goal bias \cite{lavalle1998rrt,lavalle2001randomized}}
\label{alg:ch16-rrt}
\KwIn{$\state_{\mathrm{init}}$, $\state_{\mathrm{goal}}$, goal radius $r_{\mathrm{goal}}$, step size $\eta$, goal bias $p_{\mathrm{goal}}$, iteration budget $N$}
\KwOut{a collision-free polygonal path from $\state_{\mathrm{init}}$ into $X_{\mathrm{goal}}$, or \emph{failure}}
\Fn{\RrtGrow{$\state_{\mathrm{init}}$, $\state_{\mathrm{goal}}$, $N$}}{
  $V\gets\set{\state_{\mathrm{init}}}$; $E\gets\emptyset$\label{alg:ch16-rrt:init}\;
  \For{$i\gets1$ \KwTo $N$}{
    $\state_{\mathrm{rand}}\gets$ \RrtSampleBiased{}\label{alg:ch16-rrt:sample}\tcp*{$\state_{\mathrm{goal}}$ with probability $p_{\mathrm{goal}}$, else uniform in $\mathcal{X}$}
    $\state_{\mathrm{near}}\gets\Nearest(V,\state_{\mathrm{rand}})$\label{alg:ch16-rrt:nearest}\;
    $\state_{\mathrm{new}}\gets\Steer(\state_{\mathrm{near}},\state_{\mathrm{rand}})$\label{alg:ch16-rrt:steer}\tcp*{at most $\eta$ from $\state_{\mathrm{near}}$, \cref{eq:ch16-steer}}
    \If{\RrtSegmentFree{$\state_{\mathrm{near}}$, $\state_{\mathrm{new}}$}\label{alg:ch16-rrt:free}}{
      $V\gets V\cup\set{\state_{\mathrm{new}}}$; $E\gets E\cup\set{(\state_{\mathrm{near}},\state_{\mathrm{new}})}$\label{alg:ch16-rrt:add}\;
      \If{$\norm{\state_{\mathrm{new}}-\state_{\mathrm{goal}}}\le r_{\mathrm{goal}}$\label{alg:ch16-rrt:goal}}{
        \KwRet{\RrtExtractPath{$V$, $E$, $\state_{\mathrm{new}}$}}\tcp*{follow the parents back to $\state_{\mathrm{init}}$}
      }
    }
  }
  \KwRet{failure}\;
}
\end{algorithm}

\Cref{alg:ch16-rrt} is the whole algorithm. Line~\ref{alg:ch16-rrt:init}
starts the tree with the root. Each iteration draws one sample
(line~\ref{alg:ch16-rrt:sample}); the goal bias makes every twentieth sample
(for $p_{\mathrm{goal}}=0.05$) the goal itself, so that the tree is pulled
toward the goal once in a while instead of waiting for a uniform sample to
land in the small goal region. Line~\ref{alg:ch16-rrt:nearest} is the
expensive step, a search over all vertices (\cref{sec:ch16-nearest}).
Line~\ref{alg:ch16-rrt:steer} limits the new edge to length $\eta$: without
the limit a far-away sample would produce a long edge that is almost always
blocked. Line~\ref{alg:ch16-rrt:free} is the only place where the obstacles
enter, and it tests the whole segment, not just the new vertex. If the
segment is free, the new vertex joins the tree with $\state_{\mathrm{near}}$ as
its parent (line~\ref{alg:ch16-rrt:add}); if not, the sample is simply
forgotten. The goal test of line~\ref{alg:ch16-rrt:goal} asks whether the new
vertex lies in the goal ball: exact equality with $\state_{\mathrm{goal}}$ never
happens with continuous samples, which is why a goal region is part of the
problem statement. The implementation of \cref{sec:ch16-implementation}
finally appends $\state_{\mathrm{goal}}$ itself when the segment from the entry
vertex to it is free, so that the path ends exactly at the goal.

Kuffner and LaValle \cite{kuffner2000rrtconnect} package
lines~\ref{alg:ch16-rrt:nearest}--\ref{alg:ch16-rrt:add} as one operation
$\textsc{Extend}(\mathcal{T},\state)$ with three outcomes: \textsc{Trapped} if the segment
was blocked, \textsc{Advanced} if a new vertex was added short of $\state$, and
\textsc{Reached} if the new vertex is $\state$ itself, which happens exactly when
$\norm{\state-\state_{\mathrm{near}}}\le\eta$. \Cref{alg:ch16-connect} is written in
these terms.

\paragraph{The Voronoi bias, quantified.}
With $p_{\mathrm{goal}}=0$ the sample is uniform on $\mathcal{X}$, so the
probability that vertex $\vect{v}$ is chosen by line~\ref{alg:ch16-rrt:nearest} is
\begin{equation}
\Prob\bigl(\state_{\mathrm{near}}=\vect{v}\bigr)=\frac{\mathrm{vol}\bigl(\mathrm{Vor}(\vect{v})\cap\mathcal{X}\bigr)}{\mathrm{vol}(\mathcal{X})},
\qquad \mathrm{Vor}(\vect{v})=\set{\state : \norm{\state-\vect{v}}\le\norm{\state-\vect{w}}\ \text{for all } \vect{w}\in V},
\label{eq:ch16-voronoi}
\end{equation}
the relative volume of its Voronoi region. LaValle \cite{lavalle1998rrt} observed that
the distribution of the vertices converges to the sampling distribution, so
a uniformly sampled \rrt eventually covers $\Xfree$ uniformly. Goal samples
always extend the vertex nearest to the goal, which hurts when that vertex
sits against an obstacle (\cref{sec:ch16-pitfalls}).

% ---------------------------------------------------------------------
\section{A worked example}
\label{sec:ch16-example}

\begin{example}[\rrt on a map with two walls]
\label{ex:ch16-map}
The configuration space is the square $\mathcal{X}=[0,10]^{2}$ with two
rectangular obstacles, $[3,5]\times[0,6]$ and $[6,8]\times[4,10]$, so that
the only way from the bottom-left to the top-right corner passes between
them. The robot is a point ($r=0$), the start is
$\state_{\mathrm{init}}=(1,1)$, the goal $\state_{\mathrm{goal}}=(9,9)$ with
$r_{\mathrm{goal}}=0.5$, the step size $\eta=0.5$, the goal bias
$p_{\mathrm{goal}}=0.05$ and the check resolution $\delta=0.05$. The random
generator is seeded with $1$. The instance is \code{worked\_example\_world()}
in \code{ch16\_rrt.py}, and every number below is printed by its self-test or
by the figure script \code{gen\_ch16\_tree.py}.
\end{example}

\begin{table}[tb]
\centering\small
\caption{Trace of the first ten iterations of \cref{alg:ch16-rrt} on
\cref{ex:ch16-map}. Vertices are numbered in the order of insertion, vertex~0
being the start. Iteration~4 drew the goal (goal bias); from iteration~7 on
every step is rejected by the wall $[3,5]\times[0,6]$, as the text explains.}
\label{tab:ch16-trace}
\begin{tabular}{@{}rccccl@{}}
\toprule
Iteration & $\state_{\mathrm{rand}}$ & Nearest vertex & $\state_{\mathrm{near}}$ & $\state_{\mathrm{new}}$ & Segment free? \\
\midrule
1  & $(9.50, 1.44)$ & 0 & $(1.00, 1.00)$ & $(1.50, 1.03)$ & yes, vertex 1 added\\
2  & $(3.12, 4.23)$ & 1 & $(1.50, 1.03)$ & $(1.72, 1.47)$ & yes, vertex 2\\
3  & $(4.09, 5.50)$ & 2 & $(1.72, 1.47)$ & $(1.98, 1.90)$ & yes, vertex 3\\
4  & $(9.00, 9.00)$ goal & 3 & $(1.98, 1.90)$ & $(2.33, 2.26)$ & yes, vertex 4\\
5  & $(5.38, 3.30)$ & 4 & $(2.33, 2.26)$ & $(2.80, 2.42)$ & yes, vertex 5\\
6  & $(3.03, 4.53)$ & 5 & $(2.80, 2.42)$ & $(2.86, 2.92)$ & yes, vertex 6\\
7  & $(4.03, 2.03)$ & 5 & $(2.80, 2.42)$ & $(3.28, 2.27)$ & no, enters the wall\\
8  & $(7.50, 2.80)$ & 6 & $(2.86, 2.92)$ & $(3.36, 2.90)$ & no\\
9  & $(9.81, 9.62)$ & 6 & $(2.86, 2.92)$ & $(3.22, 3.26)$ & no\\
10 & $(5.41, 2.77)$ & 6 & $(2.86, 2.92)$ & $(3.36, 2.89)$ & no\\
\bottomrule
\end{tabular}
\end{table}

\Cref{tab:ch16-trace} traces the first ten iterations. The first sample lies
far to the right, so \textsc{Steer} produces a step of exactly $\eta=0.5$
from the start. The next samples lie toward the north-east, and the tree
grows as a single chain of six vertices in that direction, exactly what the
Voronoi bias predicts for a tree with one branch: its tip owns almost the
whole map. At iteration~7 the chain has reached $x\approx2.9$ and the wall
at $x=3$ blocks every step toward the right half of the map. The tree is now
stuck in an unproductive regime: samples on the right, which are most of
them, keep selecting the two vertices closest to the wall, and
\textsc{CollisionFree} keeps rejecting the step; only samples on the left of
the wall, probability about $0.3$, add vertices. This is why the first panel
of \cref{fig:ch16-snapshots} shows only $17$ vertices after $50$
iterations: $16$ of the $50$ extensions succeeded and $34$ were rejected.

\begin{figure}[tb]
  \centering
  \inputfigure{ch16/tree-snapshots}
  \caption{The tree of \cref{ex:ch16-map} after (a)~50, (b)~200 and (c)~1000
  iterations, grown past the first solution for the last panel. The dashed
  circle is the goal region of radius $r_{\mathrm{goal}}=0.5$. In~(a) the
  tree has hit the wall $x=3$ and grows only on the left; in~(b) it has found
  the gap between the two walls; in~(c) the whole free space is covered and
  the blue path, found at iteration~402, is the chain of parents from the
  vertex that first entered the goal region. The path has $50$ vertices,
  length $23.85$, and visibly wanders: the shortest path has length $16.72$.}
  \label{fig:ch16-snapshots}
\end{figure}

\Cref{fig:ch16-snapshots} shows the tree after $50$, $200$ and $1000$
iterations. By iteration~$200$ the tree has $103$ vertices and has found the
gap between the walls; the goal region is entered at iteration~$402$, when
the tree has $245$ vertices, by a vertex at $(9.14,8.60)$ at distance
$0.42<r_{\mathrm{goal}}$ from the goal. The final segment to
$\state_{\mathrm{goal}}$ is free and is appended, which gives a path of $50$
vertices and length $23.85$. Of the $402$ iterations, $158$ were rejected by
the collision check, and the checks tested $3\,965$ points (a counter the self-test prints), about ten per
iteration, as expected for $\eta/\delta=10$. The path is far from optimal:
the shortest path in $\Xfree$ hugs the corners $(3,6)$, $(5,6)$, $(6,4)$ and
$(8,4)$ and has length $16.72$ in the limit of zero clearance, and $17.20$ when
the corners are rounded off by the clearance $0.075$ of the collision
checker; the self-test prints both. The \rrt path is $43\,\%$ longer than
the zero-clearance optimum and $39\,\%$ longer than the $17.20$ reference. \Cref{sec:ch16-properties} explains why this is not bad luck, and
\cref{sec:ch16-smoothing} how most of the excess is removed afterward.

% ---------------------------------------------------------------------
\section{Properties: completeness, optimality, complexity}
\label{sec:ch16-properties}

\subsection{Probabilistic completeness}
\label{sec:ch16-completeness}

\rrt is a randomized algorithm, and its guarantee is a statement about
probabilities. A planner is \textbf{probabilistically
complete}\index{probabilistic completeness}\index{RRT!probabilistic completeness}
if, whenever a solution with positive clearance exists -- a \emph{robustly
feasible} query, the hypothesis made precise in \cref{thm:ch16-complete} --
the probability that it has found one after $n$ iterations tends to $1$ as
$n\to\infty$. This is weaker than the
completeness of \astar on a finite graph, and it says nothing when no
solution exists: \rrt then spends its budget $N$ and reports failure,
without knowing whether the query was infeasible or merely hard.

\begin{theorem}[Probabilistic completeness of \rrt; LaValle and Kuffner 2001]
\label{thm:ch16-complete}
Let $\mathcal{X}$ be a bounded box, $\eta>0$, $0\le p_{\mathrm{goal}}<1$ and
$r_{\mathrm{goal}}>0$. Suppose there is a path $\sigma$ of finite length from
$\state_{\mathrm{init}}$ to $\state_{\mathrm{goal}}$ with \textbf{clearance}\index{clearance}
$\delta_{c}>0$: every segment that lies within distance $\delta_{c}$ of $\sigma$
is accepted by \textup{\textsc{CollisionFree}}. Then the probability that
\cref{alg:ch16-rrt} has not found a solution after $n$ iterations is at most
$m/(np)$ for constants $m$ and $p>0$ that depend on $\sigma$, $\eta$,
$\delta_{c}$, $r_{\mathrm{goal}}$ and $\mathcal{X}$ but not on $n$; in
particular it tends to $0$ as $n\to\infty$.
\end{theorem}
\begin{proof}
Let $\nu=\min(\eta/3,\ \delta_{c}/4,\ r_{\mathrm{goal}})$ and choose points
$\state_0=\state_{\mathrm{init}},\state_1,\dots,\state_m=\state_{\mathrm{goal}}$ along $\sigma$ with
$\norm{\state_{k+1}-\state_k}\le\nu$ for all $k$; a path of finite length admits such a
sequence with $m\le\lceil\mathrm{length}(\sigma)/\nu\rceil+1$. Finite length
is no restriction: the image of a path with clearance $\delta_{c}$ is
compact, so finitely many of the balls of radius $\delta_{c}/2$ centered on
its points cover it, and joining the centers of consecutive overlapping
balls gives a polygonal path of finite length inside the same tube, with
clearance at least $\delta_{c}/2$. Let
$B_k$ be the open ball of radius $\nu$ around $\state_k$; since $B_k$ lies within
$\delta_{c}$ of $\sigma$ it is contained in $\Xfree\subseteq\mathcal{X}$, so a
uniform sample lands in it with probability
$p=(1-p_{\mathrm{goal}})\,\mathrm{vol}(B_k)/\mathrm{vol}(\mathcal{X})>0$,
the same for every $k$.

\emph{Claim: if the tree has a vertex $\vect{v}\in B_k$ and the current sample
satisfies $\state_{\mathrm{rand}}\in B_{k+1}$, then the iteration adds a vertex in
$B_{k+1}$.} Let $\vect{v}'=\Nearest(V,\state_{\mathrm{rand}})$. Then
$\norm{\vect{v}'-\state_{\mathrm{rand}}}\le\norm{\vect{v}-\state_{\mathrm{rand}}}
\le\norm{\vect{v}-\state_k}+\norm{\state_k-\state_{k+1}}+\norm{\state_{k+1}-\state_{\mathrm{rand}}}<3\nu\le\eta$,
so \textsc{Steer} returns $\state_{\mathrm{new}}=\state_{\mathrm{rand}}$ itself. Both
endpoints of the segment $[\vect{v}',\state_{\mathrm{rand}}]$ lie within
$3\nu+\nu=4\nu\le\delta_{c}$ of $\state_{k+1}\in\sigma$, and a ball is convex, so the
whole segment lies within $\delta_{c}$ of $\sigma$ and \textsc{CollisionFree}
accepts it. Hence $\state_{\mathrm{rand}}\in B_{k+1}$ becomes a vertex.

Let $K_n$ be the largest $k$ such that the tree after $n$ iterations has a
vertex in $B_k$; $K_0\ge0$ because $\state_0\in B_0$ is the root. By the claim,
whenever $K_n=k<m$ and the next sample lands in $B_{k+1}$, $K_{n+1}\ge k+1$.
Samples are independent of the past and each lands in the required ball with
probability at least $p$, so the number of iterations $T$ until $K_n=m$ is
stochastically at most a sum of $m$ independent geometric variables with
success probability $p$, whence $\E[T]\le m/p$ and, by Markov's inequality,
$\Prob(T>n)\le m/(np)$. When $K_n=m$ the tree has a vertex within
$\nu\le r_{\mathrm{goal}}$ of $\state_{\mathrm{goal}}$, that is, in
$X_{\mathrm{goal}}$, and line~\ref{alg:ch16-rrt:goal} of \cref{alg:ch16-rrt}
has returned a path.
\qedhere
\end{proof}

Three things in the proof deserve attention.
\begin{itemize}
  \item \emph{The role of $\eta$.} The step size only enters through $\nu\le\eta/3$:
        a smaller $\eta$ means smaller balls, a smaller $p$ and more
        iterations, but any $\eta>0$ gives completeness. A larger $\eta$ does
        not hurt this bound, because \textsc{Steer} moves \emph{at most} $\eta$ and
        reaches nearby samples exactly.
  \item \emph{The role of the goal bias.} Only $p_{\mathrm{goal}}<1$ is needed;
        goal samples are simply wasted in the argument. With
        $p_{\mathrm{goal}}=1$ the algorithm is deterministic and fails on any
        map where the straight line from start to goal is blocked
        (\cref{exr:ch16-completeness}).
  \item \emph{Narrow passages.} The ball volume is proportional to $\nu^{d}$, so
        $p\approx(1-p_{\mathrm{goal}})\,\zeta_d\,\nu^{d}/\mathrm{vol}(\mathcal{X})$
        with $\zeta_d$ the volume of the unit ball. A passage of width $w$
        forces $\delta_{c}\le w/2$ and $\nu\le w/8$: the expected number of
        iterations grows like $(1/w)^{d}$. This \textbf{narrow-passage
        problem}\index{narrow-passage problem} is the known weakness of every
        uniform sampler and one motivation for the bidirectional search of
        \cref{sec:ch16-connect}. In \cref{ex:ch16-map} the gap is one unit
        wide on a ten-unit map and is found within a few hundred iterations.
\end{itemize}
The bound $m/(np)$ is loose; \cref{exr:ch16-completeness} sharpens it to an
exponential decay in $n$. The original argument is in
\textcite{lavalle2001randomized}; \textcite{kleinbort2019completeness} give a
modern, careful proof that also covers the kinodynamic variant of
\cref{sec:ch16-kinodynamic}, and \cref{ch:ch17} reuses the ball construction
for \rrtstar.

\subsection{Not optimal, and not even eventually optimal}
\label{sec:ch16-suboptimal}

\Cref{fig:ch16-snapshots} already showed a path $43\,\%$ longer than
necessary. \Cref{fig:ch16-jagged} shows the worst of $100$ runs of
\cref{ex:ch16-map} with different seeds: its path has length $28.35$,
$65\,\%$ above the $17.20$ reference. Over the $100$ runs the lengths range
from $20.55$ to $28.35$, with mean $23.87$ and standard deviation $1.84$.
Two features of \cref{alg:ch16-rrt}
are responsible. The parent of a vertex is fixed forever at insertion, so a
vertex reached early by a detour keeps its detour, and every later vertex
behind it inherits it. And the path zigzags at the scale of $\eta$, because
each edge points toward a random sample, not toward the goal.

\begin{figure}[tb]
  \centering
  \inputfigure{ch16/jagged}
  \caption{\rrt is not optimal. The worst of $100$ runs of
  \cref{ex:ch16-map} (seed~14, solution at iteration~746): the blue \rrt path
  has length $28.35$, the dashed orange shortest path (clearance $0.075$)
  $17.20$. The path wanders because the parent of every vertex is fixed at
  insertion and because every edge points toward a random sample. The gray
  tree shows that the tree \emph{had} vertices close to the shortest path;
  they were simply not on the chain of parents that reached the goal first.}
  \label{fig:ch16-jagged}
\end{figure}

Could one let the tree grow and wait for a better path to appear? No.
Karaman and Frazzoli \cite{karaman2011sampling} proved that this never
happens.

\begin{theorem}[\rrt converges to a suboptimal path; Karaman and Frazzoli 2011]
\label{thm:ch16-suboptimal}
Let $c^{*}$ be the length of a shortest solution and let $Y_n$ be the length
of the shortest path from $\state_{\mathrm{init}}$ into $X_{\mathrm{goal}}$
contained in the tree after $n$ iterations of \cref{alg:ch16-rrt}
(continued past the first solution). Assume, with
\textcite{karaman2011sampling}, that $X_{\mathrm{goal}}$ has non-empty
interior, that some optimal path has \emph{weak clearance} (it is the limit
of paths of positive clearance) and that $\eta$ is held fixed. Then, for
$d\ge2$,
\[
\Prob\Bigl(\lim_{n\to\infty}Y_n=c^{*}\Bigr)=0 .
\]
The best path in the tree converges almost surely to a value strictly larger
than the optimum.
\end{theorem}

The proof analyses the vertices nearest to the start along an optimal path
and shows that, with probability one, the tree's first edges in that
direction are not aligned with the optimal path, and no later sample can
repair them since edges are never changed. The remedy is to let the tree
\emph{rewire}: \rrtstar (\cref{ch:ch17}) reconsiders the parent of every new
vertex among its neighbors within a shrinking radius and reconnects existing
vertices through it when that is cheaper, which restores optimality in the
limit at the price of more collision checks. For \rrt itself the practical
remedy is post-processing (\cref{sec:ch16-smoothing}), which removes most of
the excess length but cannot carry the path to the other side of an obstacle.

\subsection{Complexity}
\label{sec:ch16-complexity}

Let $n$ be the number of vertices. One iteration costs one nearest-neighbor
query, $\bigO{n}$ with a linear scan and $\bigO{\log n}$ on average with a
kd-tree (\cref{sec:ch16-nearest}), plus one collision check of
$\lceil\eta/\delta\rceil+1$ point tests, each $\bigO{n_{\mathrm{obs}}}$ in
the number of obstacles. After $N$ iterations the total is $\bigO{N^{2}}$
with the scan and $\bigO{N\log N}$ with a kd-tree, plus
$\bigO{N\,n_{\mathrm{obs}}\,\eta/\delta}$ for the checks, which dominate in
cluttered scenes. The memory is $\bigO{N}$ vertices with $d$ coordinates and
one parent index each, whatever the size or dimension of $\mathcal{X}$,
which is the whole point of \cref{tab:ch16-gridsizes}. What no bound tells
you is \emph{how many} iterations are needed; that depends on the narrowest
passage the path must cross.

% ---------------------------------------------------------------------
\section{Variants and extensions}
\label{sec:ch16-variants}

\subsection{\rrt-Connect: two trees and a greedy connection}
\label{sec:ch16-connect}

A single tree rooted at the start has to discover the goal region by chance
or by goal bias. \textcite{kuffner2000rrtconnect} grow two trees, one from
the start and one from the goal, and add a greedy step: whenever one tree
gains a new vertex, the other tree tries to reach that vertex by repeating
\textsc{Extend} as long as it makes progress. \Cref{alg:ch16-connect} gives
both ingredients.\index{RRT-Connect}\index{RRT!RRT-Connect}

\begin{algorithm}[htb]
\caption{\rrt-Connect \cite{kuffner2000rrtconnect}}
\label{alg:ch16-connect}
\KwIn{$\state_{\mathrm{init}}$, $\state_{\mathrm{goal}}$, step size $\eta$, budget $N$}
\KwOut{a collision-free path from $\state_{\mathrm{init}}$ to $\state_{\mathrm{goal}}$, or \emph{failure}}
\Fn{\RrtExtendTree{$\mathcal{T}$, $\state$}}{
  $\state_{\mathrm{near}}\gets\Nearest(V_{\mathcal{T}},\state)$; $\state_{\mathrm{new}}\gets\Steer(\state_{\mathrm{near}},\state)$\;
  \If{\RrtSegmentFree{$\state_{\mathrm{near}}$, $\state_{\mathrm{new}}$}}{
    add $\state_{\mathrm{new}}$ to $\mathcal{T}$ with parent $\state_{\mathrm{near}}$\;
    \lIf{$\state_{\mathrm{new}}=\state$}{\KwRet{\textup{\textsc{Reached}}}}
    \KwRet{\textup{\textsc{Advanced}}}\;
  }
  \KwRet{\textup{\textsc{Trapped}}}\;
}
\Fn{\RrtConnectGreedy{$\mathcal{T}$, $\state$}}{
  \Repeat{$s\neq\textup{\textsc{Advanced}}$}{$s\gets$ \RrtExtendTree{$\mathcal{T}$, $\state$}\label{alg:ch16-connect:greedy}\;}
  \KwRet{$s$}\;
}
\Fn{\RrtConnectPlan{$\state_{\mathrm{init}}$, $\state_{\mathrm{goal}}$, $N$}}{
  $\mathcal{T}_a\gets$ tree with root $\state_{\mathrm{init}}$; $\mathcal{T}_b\gets$ tree with root $\state_{\mathrm{goal}}$\;
  \For{$i\gets1$ \KwTo $N$}{
    $\state_{\mathrm{rand}}\gets$ uniform sample in $\mathcal{X}$\tcp*{no goal bias}
    \If{\RrtExtendTree{$\mathcal{T}_a$, $\state_{\mathrm{rand}}$} $\neq\textup{\textsc{Trapped}}$\label{alg:ch16-connect:extend}}{
      $\state_{\mathrm{new}}\gets$ the vertex just added to $\mathcal{T}_a$\;
      \If{\RrtConnectGreedy{$\mathcal{T}_b$, $\state_{\mathrm{new}}$} $=\textup{\textsc{Reached}}$\label{alg:ch16-connect:test}}{
        \KwRet{the path root$(\mathcal{T}_a)\to \state_{\mathrm{new}}$ in $\mathcal{T}_a$ joined with $\state_{\mathrm{new}}\to$ root$(\mathcal{T}_b)$ in $\mathcal{T}_b$, reversed if $\mathcal{T}_a$ is the goal tree}\;
      }
    }
    \RrtSwap{$\mathcal{T}_a$, $\mathcal{T}_b$}\label{alg:ch16-connect:swap}\;
  }
  \KwRet{failure}\;
}
\end{algorithm}

Line~\ref{alg:ch16-connect:extend} extends the current tree $\mathcal{T}_a$ toward a
uniform sample exactly as \rrt does. If that produced a vertex, the other
tree $\mathcal{T}_b$ tries to connect to it (line~\ref{alg:ch16-connect:test}): the
\textsc{Connect} heuristic\index{Connect heuristic} of
line~\ref{alg:ch16-connect:greedy} repeats \textsc{Extend} toward the
\emph{same} target, adding one edge of length $\eta$ per repetition, until
the target is reached or an obstacle is hit. In open space this is a straight
run of edges across the map, and this greedy extension is what makes the
method fast. The trees swap roles every iteration
(line~\ref{alg:ch16-connect:swap}), which keeps them balanced; no goal bias
is needed, because the goal is the root of the second tree. When
\textsc{Connect} returns \textsc{Reached}, $\state_{\mathrm{new}}$ belongs to both
trees and the two root-to-vertex chains form the path, one of them reversed;
it is collision-free because every edge of either tree is.
\Cref{fig:ch16-connect} shows the two trees on \cref{ex:ch16-map}: with
seed~1 they meet at iteration~$164$ with $68+29$ vertices, against $402$
iterations and $245$ vertices for the single tree, and the path has length
$25.11$, no better than before.

\begin{figure}[tb]
  \centering
  \inputfigure{ch16/connect}
  \caption{\rrt-Connect on \cref{ex:ch16-map} (seed~1). The blue tree grows
  from the start, the orange tree from the goal; the trees meet at
  iteration~164 with $68+29$ vertices. The long straight runs of edges are the
  work of \textsc{Connect}: each one is a sequence of \textsc{Advanced}
  results toward one target vertex of the other tree. The path (thick blue)
  is as jagged as that of the single tree.}
  \label{fig:ch16-connect}
\end{figure}

\begin{figure}[tb]
  \centering
  \inputfigure{ch16/iterations}
  \caption{Fraction of $100$ runs of \cref{ex:ch16-map} (seeds $0$ to $99$)
  that have found a path within a given number of iterations, for \rrt with
  goal bias $0$, $0.05$ and $0.5$ and for \rrt-Connect. The empirical
  distribution of \rrt-Connect lies far to the left; a goal bias of $0.5$ is
  worse than no bias at all on this map.}
  \label{fig:ch16-iterations}
\end{figure}

\begin{table}[!t]
\centering\small
\caption{Iterations to the first solution over $100$ seeds on
\cref{ex:ch16-map}, budget $3\,000$ iterations, $\eta=0.5$; all runs succeed.
Vertices and path lengths are averages, the lengths before any smoothing.
\rrt-Connect needs $2.4$ times fewer iterations than \rrt with the default
bias and stores $2.7$ times fewer vertices, but its paths are no shorter.}
\label{tab:ch16-connect}
\begin{tabular}{@{}lrrrr@{}}
\toprule
Planner & Median iterations & Mean iterations & Vertices & Path length\\
\midrule
\rrt, $p_{\mathrm{goal}}=0$    & 518 & 537 & 352 & $24.31\pm2.00$\\
\rrt, $p_{\mathrm{goal}}=0.05$ & 425 & 453 & 275 & $23.87\pm1.84$\\
\rrt, $p_{\mathrm{goal}}=0.5$  & 745 & 783 & 245 & $23.61\pm1.86$\\
\rrt-Connect                   & 178 & 192 & 101 & $23.69\pm1.98$\\
\bottomrule
\end{tabular}
\end{table}

One run proves nothing about a randomized algorithm, so
\cref{fig:ch16-iterations} and \cref{tab:ch16-connect} report $100$ runs with
seeds $0$ to $99$. \rrt-Connect solves every instance within $375$
iterations; by iteration $500$ plain \rrt with the default goal bias has
solved $73\,\%$, and its median is $425$ iterations against $178$. The self-test of
\code{ch16\_rrt.py} asserts this ordering of medians and means over $40$
seeds. The table also isolates the goal bias: a small bias helps
($518\to425$ median iterations), a large one hurts badly ($745$), for the
reason given in \cref{sec:ch16-pitfalls}. The tree
with $p_{\mathrm{goal}}=0.5$ has \emph{fewer} vertices after \emph{more}
iterations, the signature of wasted, trapped extensions.

\subsection{Kinodynamic \rrt}
\label{sec:ch16-kinodynamic}


The paths of \cref{alg:ch16-rrt} are polygons, and a drone cannot fly a
polygon at speed: it has bounded acceleration, so its velocity is part of its
state. \textcite{lavalle2001randomized} extended \rrt to systems with
dynamics; the version below uses the double integrator of \cref{ch:ch02},
with state $\state=(\pos,\vel)\in\R^{2d}$, control $\acc$ with
$\norm{\acc}\le a_{\max}$ held constant over a step $\dt$, and speed limit
$\norm{\vel}\le v_{\max}$.\index{kinodynamic planning}\index{RRT!kinodynamic}\index{double integrator!kinodynamic RRT}
Two primitives change. \textsc{Steer} becomes forward simulation: the
planner cannot move ``toward'' a sampled state along a straight line,
because the straight line is not a feasible motion. Instead it tries $m$
random controls, integrates each for one step, keeps the result that ends
closest to the sample, and stores the control on the edge. \textsc{Nearest}
becomes a question rather than a formula: the Euclidean distance between two
states is not the time or effort needed to move between them, since a state
whose velocity points away from the sample is ``far'' even when its position
is close. The exact answer, the cost of the optimal motion between two
states, requires solving a two-point boundary value problem for every
candidate vertex, far too expensive inside a nearest-neighbor query. The
usual compromise is the weighted metric
\begin{equation}
d_{w}(\state,\state')=\sqrt{\norm{\pos-\pos'}^{2}+w_{\mathrm{v}}^{2}\norm{\vel-\vel'}^{2}},
\label{eq:ch16-metric}
\end{equation}
with a weight $w_{\mathrm{v}}$ (units of time) that says how much a velocity
mismatch counts against a position mismatch. \Cref{alg:ch16-kino} is the
result; line~\ref{alg:ch16-kino:integrate} is the exact zero-order-hold
update of \cref{ch:ch02}. The arc flown during one step is a parabola, and
the collision check samples it at spacing $\delta$, using the fact that the
speed along a constant-acceleration arc is largest at one of its ends.

\begin{algorithm}[!htb]
\caption{Kinodynamic \rrt for the double integrator \cite{lavalle2001randomized}}
\label{alg:ch16-kino}
\KwIn{$\state_{\mathrm{init}}=(\pos_{\mathrm{init}},\vect{0})$, $\pos_{\mathrm{goal}}$, $r_{\mathrm{goal}}$, $v_{\max}$, $a_{\max}$, $\dt$, controls per extension $m$, goal bias $p_{\mathrm{goal}}$, velocity weight $w_{\mathrm{v}}$, budget $N$}
\KwOut{a sequence of states and constant accelerations that obeys the dynamics, or \emph{failure}}
\Fn{\RrtKinoGrow{$\state_{\mathrm{init}}$, $\pos_{\mathrm{goal}}$, $N$}}{
  $V\gets\set{\state_{\mathrm{init}}}$\;
  \For{$i\gets1$ \KwTo $N$}{
    \lIf{$\mathrm{rand}()<p_{\mathrm{goal}}$}{$\state_{\mathrm{rand}}\gets(\pos_{\mathrm{goal}},\vect{0})$}
    \lElse{$\state_{\mathrm{rand}}\gets(\pos,\vel)$ with $\pos$ uniform in $\mathcal{X}$ and $\vel$ uniform in the ball $\norm{\vel}\le v_{\max}$}
    $\state_{\mathrm{near}}\gets\argmin_{\state\in V}d_{w}(\state,\state_{\mathrm{rand}})$\label{alg:ch16-kino:nearest}\tcp*{weighted metric \cref{eq:ch16-metric}}
    $\mathrm{best}\gets\KwNil$; $\acc_{\mathrm{best}}\gets\KwNil$\tcp*{convention $d_{w}(\KwNil,\cdot)=+\infty$}
    \For{$j\gets1$ \KwTo $m$}{
      $\acc_j\gets$ uniform in the ball $\norm{\acc}\le a_{\max}$\;
      $\state_j\gets$ \RrtIntegrate{$\state_{\mathrm{near}}$, $\acc_j$, $\dt$}\label{alg:ch16-kino:integrate}\tcp*{$\pos+\dt\,\vel+\tfrac12\dt^{2}\acc_j$, $\vel+\dt\,\acc_j$}
      \If{$\norm{\vel_j}\le v_{\max}$ \KwAnd the arc from $\state_{\mathrm{near}}$ to $\state_j$ is collision-free \KwAnd $d_{w}(\state_j,\state_{\mathrm{rand}})<d_{w}(\mathrm{best},\state_{\mathrm{rand}})$}{
        $\mathrm{best}\gets\state_j$; $\acc_{\mathrm{best}}\gets\acc_j$\;
      }
    }
    \If{$\mathrm{best}\neq\KwNil$}{
      add $\mathrm{best}$ to $V$ with parent $\state_{\mathrm{near}}$ and edge label $\acc_{\mathrm{best}}$\;
      \lIf{$\norm{\pos_{\mathrm{best}}-\pos_{\mathrm{goal}}}\le r_{\mathrm{goal}}$}{\KwRet{states and controls from the root to $\mathrm{best}$}}
    }
  }
  \KwRet{failure}\;
}
\end{algorithm}


The implementation \code{kinodynamic\_rrt} in \code{ch16\_rrt.py} solves
\cref{ex:ch16-map} for a drone with $v_{\max}=2$, $a_{\max}=2$, $\dt=0.5$,
$m=8$ controls per extension, $w_{\mathrm{v}}=0.5$ and seed~2: it finds a
trajectory at iteration~$680$ with $512$ vertices, $41$ steps or $20.5$~s of
flight, arriving at speed $0.54$. The self-test verifies that consecutive
states satisfy $\state_{k+1}=\mat{A}\state_k+\mat{B}\acc_k$ with the matrices
of \cref{ch:ch02} and that no control or velocity exceeds its limit. Three
lessons carry over to every planner with dynamics: the tree needs more
vertices ($512$ against $245$) because velocities add dimensions; the
solution is a trajectory that can be flown as it is; and the metric of
line~\ref{alg:ch16-kino:nearest} strongly affects the growth of the tree.
Kinodynamic planning has a large literature
\cite{lavalle2006planning}; in the hybrid architecture of \cref{ch:ch24}
the geometric \rrt with post-processing, followed by the controller of
\cref{ch:ch21}, is usually the simpler choice.

\subsection{Post-processing: shortcutting and smoothing}
\label{sec:ch16-smoothing}


A raw \rrt path is safe but ugly. \textbf{Shortcutting}\index{shortcutting}\index{path smoothing}
repairs most of the damage with the same collision checker: pick two points
on the path, and if the straight segment between them is free, replace the
part of the path in between by the segment. By the triangle inequality the
length can only decrease. \Cref{alg:ch16-shortcut} gives the randomized
version of \textcite{geraerts2007creating} and a greedy variant.

\begin{algorithm}[!htb]
\caption{Path post-processing: random shortcutting and greedy pruning}
\label{alg:ch16-shortcut}
\KwIn{a collision-free polygonal path $\sigma=(\state_0,\dots,\state_k)$, number of attempts $K$}
\KwOut{a collision-free path with the same endpoints that is no longer than $\sigma$}
\Fn{\RrtShortcutPath{$\sigma$, $K$}}{
  \For{$i\gets1$ \KwTo $K$}{
    draw $s_1<s_2$ uniformly from $[0,\mathrm{length}(\sigma)]$ and let $\vect{a}=\sigma(s_1)$, $\vect{b}=\sigma(s_2)$ be the points at these arc lengths\;
    \If{$\vect{a}$ and $\vect{b}$ lie on different segments of $\sigma$ \KwAnd \RrtSegmentFree{$\vect{a}$, $\vect{b}$}}{
      replace the part of $\sigma$ between $\vect{a}$ and $\vect{b}$ by the segment $[\vect{a},\vect{b}]$\label{alg:ch16-shortcut:replace}\;
    }
  }
  \KwRet{$\sigma$}\;
}
\Fn{\RrtPrunePath{$\sigma$}}{
  $\sigma'\gets(\state_0)$; $i\gets0$\;
  \While{$i<k$}{
    $j\gets$ the largest index $>i$ with \RrtSegmentFree{$\state_i$, $\state_j$}\label{alg:ch16-shortcut:jump}\;
    append $\state_j$ to $\sigma'$; $i\gets j$\;
  }
  \KwRet{$\sigma'$}\;
}
\end{algorithm}


\begin{figure}[tb]
  \centering
  \inputfigure{ch16/shortcut}
  \caption{Post-processing of the path of \cref{ex:ch16-map}. Gray: the raw
  \rrt path ($50$ vertices, length $23.85$). Blue: after $100$ random
  shortcut attempts ($23$ vertices, $19.91$). Dashed orange with dots: after
  greedy pruning ($7$ vertices, $18.39$). The shortest path with clearance
  $0.075$ has length $17.20$. Both results are collision-free, and the pruned
  path hugs the obstacle corners, which is why the obstacles should be
  inflated by a safety margin before planning.}
  \label{fig:ch16-shortcut}
\end{figure}

\Cref{fig:ch16-shortcut} shows the effect on the worked example: $100$
random attempts bring the length from $23.85$ to $19.91$, and greedy pruning
to $18.39$ with only $7$ vertices, within $7\,\%$ of the shortest path with
the same clearance. The self-test asserts that neither method ever increases
the length and that the results are collision-free. Pruning hugs the
obstacle corners as tightly as the checker allows, good for length and bad
for safety, so plan with obstacles inflated by a margin beyond the drone
radius. Both methods improve the path locally: every shortcut segment must
itself be free, so a path that takes the long way around a large obstacle
keeps taking it (\cref{exr:ch16-shortcut}), and only a new search, or
\rrtstar (\cref{ch:ch17}), finds the better route. The polygonal result is
turned into a flyable trajectory by the controller of \cref{ch:ch21}, or by
fitting a smooth curve with bounded velocity and acceleration through the
waypoints, which must then be collision-checked again. In practice the
corners are rounded first, each waypoint being replaced by a short arc whose
radius respects the speed and acceleration limits of \cref{ch:ch02}; the arc
is then sampled at the resolution $\delta$ and checked again, because
cutting a corner moves the path toward the inside of the turn, which is
exactly where the obstacle that caused the corner sits. \Cref{ch:ch21} takes
the better route of optimizing the trajectory subject to the dynamics and
the obstacles instead of repairing it afterward.

\subsection{Planning in 3D}
\label{sec:ch16-3d}

\begin{figure}[!t]
  \centering
  \inputfigure{ch16/scene3d}
  \caption{\rrt in the 3D box field \code{scene\_3d()}, drawn in an oblique
  projection ($y$ recedes into the page). The light tree has $158$ vertices;
  the blue path (length $24.48$) flies over the wall on the left and around
  the pillar, and the dashed orange path is its shortcut ($20.40$). The code
  is the same as in 2D; only the bounds of $\mathcal{X}$ have three
  components.}
  \label{fig:ch16-scene3d}
\end{figure}


Nothing in \cref{alg:ch16-rrt} refers to the dimension. The code of
\cref{sec:ch16-implementation} takes the bounds of $\mathcal{X}$ as two
vectors, and the only difference between the 2D and the 3D planner is that
they have three components: \code{World([0, 0, 0], [10, 10, 10],
boxes=...)} instead of \code{World([0, 0], [10, 10], boxes=...)}, and the
same \code{rrt()} call. \Cref{fig:ch16-scene3d} shows the 3D box field
\code{scene\_3d()}, a wall to fly over, a full-height pillar, a slab to fly
under or over and three blocks, for a drone of radius $0.3$ with
$\delta=0.1$ and $\eta=0.8$. With seed~3, \rrt reaches the goal at
iteration~$321$ with $158$ vertices and a path of length $24.48$;
shortcutting brings it to $20.40$, against a straight-line distance of
$15.59$. \rrt-Connect solves the same query at iteration~$117$ with $59$
vertices. A $0.1$~m grid of this $10\times10\times10$ scene has a million
cells; \rrt has looked at a few hundred points.

\subsection{Relatives}
\label{sec:ch16-relatives}

The \textbf{probabilistic roadmap} (PRM)\index{probabilistic roadmap (PRM)}
of \textcite{kavraki1996prm} is the other founding method of sampling-based
planning: sample configurations, keep the free ones, connect each to its $k$
nearest neighbors by collision-checked segments, and answer any number of
queries by graph search (\cref{ch:ch03,ch:ch04}) on the roadmap. PRM is a
\emph{multi-query} method that spends its effort once per map; \rrt is a
\emph{single-query} method that spends it once per query and never builds a
graph of the whole space, which suits a drone whose map and queries change
every time and is the only option when the dynamics allow no exact steering.
\rrtstar and \irrtstar \cite{karaman2011sampling,gammell2014informed} keep
the tree but make it optimal in the limit; they are the subject of
\cref{ch:ch17}, together with BIT\textsuperscript{*} \cite{gammell2015bit}.
Both books \cite{lavalle2006planning,choset2005principles} survey the many
other variants.

% ---------------------------------------------------------------------
\section{Implementation notes}
\label{sec:ch16-implementation}

\subsection{Collision checking for drones}
\label{sec:ch16-collision}

The obstacles of this chapter are axis-aligned boxes and spheres, which
cover buildings, no-fly volumes and other drones well enough for planning
and have closed-form distance functions.\index{collision checking!sphere versus box}
\begin{itemize}
  \item \emph{Sphere versus box.} The distance from a point $\state$ to the box
        $[\vect{l},\vect{u}]$ is the length of the vector from $\state$ to its
        \emph{clamped} copy, $\norm{\state-\mathrm{clamp}(\state,\vect{l},\vect{u})}$,
        with $\mathrm{clamp}$ applied per coordinate; the sphere of radius $r$
        around $\state$ is free of the box if and only if this distance exceeds $r$.
        \code{World.box\_distances} computes it for all points and all boxes
        at once with NumPy broadcasting.
  \item \emph{Segment versus sphere.} A sphere obstacle of radius $R$ at
        $\vect{c}$ meets the sphere robot moving along $[\vect{a},\vect{b}]$ if and only
        if the distance from $\vect{c}$ to the segment is at most $R+r$; the
        point-to-segment distance of \cref{ch:ch02} (project $\vect{c}$ onto
        the line, clamp the parameter to $[0,1]$) answers this exactly, with
        no subdivision.
  \item \emph{Conservative inflation.} Boxes are checked by subdivision
        with the margin $r+\delta/2$ of \cref{thm:ch16-resolution}, so a path
        that passes the check keeps a clearance of at least $r$ from every
        obstacle. Any extra safety margin, for the position uncertainty of
        \cref{ch:ch18} for example, is simply added to $r$.
\end{itemize}
\Cref{lst:ch16-segment} is the segment test; the vertex test
\code{point\_free} is the special case of a segment of length zero.

\begin{lstlisting}[caption={The segment test of \code{World} (from \code{code/ch16\_rrt.py}): subdivision with the conservative margin $r+\delta/2$ for the boxes, the exact segment--sphere test for the spheres.},label={lst:ch16-segment}]
    def segment_free(self, a, b):
        """CollisionFree(a, b): subdivision check of the straight segment.

        Points spaced at most delta apart are tested against the boxes
        inflated by radius + delta/2 (conservative: every point of the
        segment is within delta/2 of a tested point, so the true clearance
        is at least radius).  Spheres use the exact segment-sphere test."""
        a = np.asarray(a, dtype=float)
        b = np.asarray(b, dtype=float)
        n = max(1, int(math.ceil(np.linalg.norm(b - a) / self.delta)))
        pts = a + np.linspace(0.0, 1.0, n + 1)[:, None] * (b - a)
        if not np.all(self.points_free(pts, self.radius + 0.5 * self.delta,
                                       spheres=False)):
            return False
        if len(self.sph_c):
            ab = b - a
            t = (self.sph_c - a) @ ab / max(float(ab @ ab), 1e-300)
            closest = a + np.clip(t, 0.0, 1.0)[:, None] * ab
            d = np.linalg.norm(self.sph_c - closest, axis=1)
            if np.any(d <= self.sph_r + self.radius):
                return False
        return True
\end{lstlisting}

\subsection{Nearest neighbors and kd-trees}
\label{sec:ch16-nearest}

The tree of \code{ch16\_rrt.py} keeps its vertices in one NumPy array and
answers \textsc{Nearest} by a vectorized linear scan: one subtraction, one
squared norm per row and an \code{argmin}. That is $\bigO{n}$ per query with
a small constant, about $0.04$~ms for $n=10^{3}$ vertices, $0.3$~ms for
$10^{4}$ and $3$~ms for $10^{5}$ in three dimensions: means of $200$ queries,
timed and printed by the self-test of \code{ch16\_rrt.py} on the machine
that produced this book, and reproducible to within a factor of two on any
other. For the trees of this chapter the scan is not the bottleneck, and it
needs no rebuilding when a vertex is added. Beyond $10^{4}$ vertices use a
\textbf{kd-tree}\index{kd-tree}\index{nearest-neighbor search} \cite{bentley1975kdtree}: a
binary tree that splits the points along one coordinate axis at each level,
so that a query descends to the leaf containing the query point and then
backtracks only into branches whose bounding region is closer than the best
distance found so far. In low dimension the expected cost is $\bigO{\log n}$
per query after an $\bigO{n\log n}$ build. Since \rrt inserts a vertex every
iteration and a balanced kd-tree does not take insertions well,
implementations rebuild it every few hundred insertions and scan the
vertices added since (\code{scipy.spatial.cKDTree} is the standard \python
choice), or use a bucket grid over $\mathcal{X}$. In high dimension the
backtracking visits most branches and a kd-tree degrades to a scan; beyond
ten or so dimensions approximate methods are used instead.

\subsection{Parameters}
\label{sec:ch16-parameters}

\begin{table}[tb]
\centering\small
\caption{The parameters of \rrt and what happens when they are badly chosen.
Values are for the map size $L$ of \cref{ex:ch16-map} ($L=10$).}
\label{tab:ch16-parameters}
\begin{tabular}{@{}lL{2.4cm}L{3.6cm}L{4.0cm}@{}}
\toprule
Parameter & Typical value & Too small & Too large\\
\midrule
Step size $\eta$ & $0.02L$ to $0.1L$ (here $0.5$) & slow growth, many vertices, long nearest-neighbor scans & long segments rejected in clutter, jagged paths; never a safety problem when the whole segment is checked\\
Goal bias $p_{\mathrm{goal}}$ & $0.05$ to $0.1$ & no pull toward the goal; the region is found by chance & greedy, trapped behind obstacles (\cref{tab:ch16-connect})\\
Goal radius $r_{\mathrm{goal}}$ & about $\eta$ & the goal region is never entered & the path ends far from the goal; always append the final segment if free\\
Check resolution $\delta$ & $\eta/10$; at most half the thinnest obstacle & slow collision checks & obstacles missed unless inflated by $\delta/2$; then over-conservative in tight spaces\\
Budget $N$ & $10^{3}$ to $10^{5}$ & failure reported on feasible queries & time wasted on infeasible queries\\
Random seed & fixed per experiment & \multicolumn{2}{L{7.6cm}}{a result from one seed is an anecdote; report statistics over seeds}\\
\bottomrule
\end{tabular}
\end{table}

\Cref{tab:ch16-parameters} lists the parameters. Two deserve a word beyond
the table: seed the generator (\code{numpy.random.default\_rng(seed)}) and
report statistics over many seeds, as in \cref{tab:ch16-connect}, never a
single run; and choose the budget $N$ from the time you can afford, since
\rrt cannot recognize an infeasible query and failure means only ``no path
found within the budget''.
\Cref{lst:ch16-rrt} is the main loop of the implementation; it matches
\cref{alg:ch16-rrt} line by line and adds the bookkeeping for the figures.
\code{extend}, \code{connect} and \code{rrt\_connect} implement
\cref{alg:ch16-connect} in the same file. The self-test runs in a few
seconds -- it prints its own runtime, three to five seconds on the machines
used for this book -- and re-checks in one place
every claim of this chapter that a program can check: the collision tests,
the worked example and its two reference lengths, the smoothing bounds, the
\rrt-Connect ordering over $40$ seeds, the 3D scene, the kinodynamic limits
and the nearest-neighbor timings quoted above.

\begin{lstlisting}[caption={The loop of \code{rrt(world, start, goal, eta, p\_goal, r\_goal, max\_iters, seed, ...)} (from \code{code/ch16\_rrt.py}). \code{Tree} stores the vertices in a growing array with parent indices; \code{steer} is \cref{eq:ch16-steer}; the goal itself is appended to the path when the last segment is free.},label={lst:ch16-rrt}]
    rng = np.random.default_rng(seed)
    start = np.asarray(start, dtype=float)
    goal = np.asarray(goal, dtype=float)
    tree = Tree(start, max_iters + 2)
    goal_index, solution_iter, snaps = -1, None, {}
    it = 0
    for it in range(1, max_iters + 1):
        q_rand = goal if rng.random() < p_goal else world.sample(rng)
        i_near = tree.nearest(q_rand)
        q_near = tree.V[i_near]
        q_new = steer(q_near, q_rand, eta)
        free = world.segment_free(q_near, q_new)
        if trace is not None:
            trace.append((it, q_rand.copy(), i_near, q_new.copy(), free))
        if free and np.linalg.norm(q_new - q_near) > 1e-12:
            i_new = tree.add(q_new, i_near)
            if goal_index < 0 and np.linalg.norm(q_new - goal) <= r_goal:
                goal_index, solution_iter = i_new, it
                if stop_at_goal:
                    break
        if it in snapshots:
            snaps[it] = tree.n
    path = None
    if goal_index >= 0:
        path = tree.path_to(goal_index)
        if (np.linalg.norm(path[-1] - goal) > 0
                and world.segment_free(path[-1], goal)):
            path = np.vstack([path, goal])
    return Result(path=path, iterations=it, solution_iter=solution_iter,
                  tree=tree, snapshots=snaps, n_vertices=tree.n)
\end{lstlisting}

\subsection{Pitfalls}
\label{sec:ch16-pitfalls}

\begin{pitfall}[Testing the vertices but not the segments]
The most common bug in a first \rrt is to test only $\state_{\mathrm{new}}$, or
the endpoints of the edge, for collision. The tree then tunnels through every
obstacle thinner than $\eta$, and the bug is invisible on maps with thick
walls. A step size that is ``too large'' is dangerous only in this sense;
with the subdivision test of \cref{thm:ch16-resolution} a long segment is
simply rejected more often. The same holds for the resolution: test points
spaced $\delta$ apart \emph{without} the inflation by $\delta/2$ miss any
wall thinner than $\delta$, at a rate that depends on the angle of the
segment; this is the trap that the self-test of \code{ch16\_rrt.py}
reproduces with the thin wall of \cref{sec:ch16-primitives}.
\end{pitfall}

\begin{pitfall}[A goal bias that is too high]
Goal bias looks like a free lunch: more samples at the goal, faster
solutions. \Cref{tab:ch16-connect} shows the opposite for
$p_{\mathrm{goal}}=0.5$: the median number of iterations rises from $425$
to $745$, above the value without any bias. Every goal sample extends the
vertex nearest to the goal, and until the tree has passed the last obstacle
that vertex sits against a wall and the extension is rejected; the tree
becomes greedy and stuck. Keep $p_{\mathrm{goal}}$ between $0.05$ and
$0.1$, or use \rrt-Connect, which needs no bias.
\end{pitfall}

\begin{pitfall}[Forgetting the goal region]
With continuous samples the tree never lands exactly on $\state_{\mathrm{goal}}$,
so a test $\state_{\mathrm{new}}=\state_{\mathrm{goal}}$ succeeds only through a goal
sample within $\eta$ of its nearest vertex, and never when
$p_{\mathrm{goal}}=0$. Define the goal as a region of radius
$r_{\mathrm{goal}}$ (\cref{def:ch16-problem}), test membership on every new
vertex, and append the final segment to $\state_{\mathrm{goal}}$ when it is free
so that the path ends where the mission requires. In the kinodynamic planner
the goal state should also allow a range of arrival velocities.
\end{pitfall}

% ---------------------------------------------------------------------
\section{Grid search or sampling?}
\label{sec:ch16-choice}

The Week~8 milestone of the study plan is to choose between graph search and
sampling-based planning; \cref{tab:ch16-comparison} collects the criteria
and its last row answers in one line. Two entries deserve a comment: only a
grid lets many agents share one discretization, which the multi-agent
planners of \cref{ch:ch07,ch:ch08,ch:ch09,ch:ch10,ch:ch11} all need, and
only sampling in the state space respects the dynamics during planning.
\Cref{exr:ch16-coding} asks you to measure the trade-off yourself.

\begin{table}[tb]
\centering\small
\caption{Grid \astar, \rrt and \rrtstar compared. ``Optimal'' for grid
\astar means optimal among grid paths, which are longer than the true
shortest path by up to $8\,\%$ on an 8-connected grid.}
\label{tab:ch16-comparison}
\begin{tabular}{@{}lL{3.3cm}L{3.3cm}L{3.3cm}@{}}
\toprule
Property & Grid \astar (\cref{ch:ch04}) & \rrt & \rrtstar (\cref{ch:ch17})\\
\midrule
Completeness & complete on the grid; blind to gaps narrower than a cell & probabilistically complete (\cref{thm:ch16-complete}) & probabilistically complete\\
Optimality & optimal on the grid & no, and never in the limit (\cref{thm:ch16-suboptimal}) & asymptotically optimal\\
Memory & the whole grid, $N^{d}$ cells & $n$ vertices & $n$ vertices\\
Anytime behavior & no; \arastar (\cref{ch:ch06}) & first solution only; improve by shortcutting & the path improves with every iteration\\
Dynamics & only via a state lattice & kinodynamic \rrt (\cref{sec:ch16-kinodynamic}) & needs an exact steering function; hard\\
Cost driven by & $N^{d}$ and the heuristic & narrow passages; nearest-neighbor and collision checks & the same plus the neighbor queries and rewiring\\
Prefer when & 2D or coarse 3D maps; many agents on one graph; exact answers & large or high-dimensional spaces; a feasible path fast & path quality matters and time allows\\
\bottomrule
\end{tabular}
\end{table}

% ---------------------------------------------------------------------
\section{Where this fits in the drone system}
\label{sec:ch16-drone}

\begin{dronebox}
In the hybrid architecture of \cref{ch:ch24} the global planner
(conflict-based search, \cbs, or its bounded-suboptimal variant \ecbs) plans
time-indexed paths for the whole swarm on a coarse grid, and the replanning
layer repairs a path with the incremental \dstarlite or a fresh \astar search
on the same grid. \rrt is the \textbf{continuous 3D planner} that steps in when the grid
is not enough: when the cells are too coarse for the local geometry (a gap
between two cranes narrower than a cell), when the volume is too large to
discretize at the accuracy the mission needs (\cref{tab:ch16-gridsizes}),
or when an unknown obstacle detected by the tracker of \cref{ch:ch18}
blocks the nominal route and a detour has to be found in three dimensions
within a fraction of a second.

\textbf{Combining a global grid plan with a local \rrt.} When a segment
between two consecutive waypoints of the grid path becomes infeasible, the
local planner runs \cref{alg:ch16-rrt} or \cref{alg:ch16-connect} from the
current position to the next reachable waypoint, with a goal region of
about one cell and the obstacles inflated by the drone radius, the tracking
error and the prediction uncertainty of \cref{ch:ch19,ch:ch20}. The result
is shortcut (\cref{alg:ch16-shortcut}), turned into a trajectory by the
controller of \cref{ch:ch21} and checked against the reservations of the
other drones; if it violates them, or if the local planner fails within its
budget, the event is escalated to a global replan, the decision logic of
\cref{ch:ch24}. In the other direction, the waypoints of the grid path can
seed the goal bias of the tree, turning the global plan into a heuristic for
the local sampler.

\textbf{What it receives and delivers.} A start, a goal region, a map of
boxes and the predicted volumes of moving obstacles come in; a
collision-free polygonal path with a known clearance, or a failure within a
known time, goes out. \rrt guarantees nothing about optimality
(\cref{thm:ch16-suboptimal}) and nothing at all when no path exists, which
is why its budget must be chosen from the available time. Week~8 of the
study plan (\cref{ch:appA}) covers this chapter and \cref{ch:ch17}; its
coding exercise is \cref{exr:ch16-coding}.
\end{dronebox}

% ---------------------------------------------------------------------
\section{Summary}
\begin{summary}
\begin{itemize}
  \item Grids cost $N^{d}$ cells; a $100\times100\times50$~m volume at
        $0.1$~m has $5\times10^{8}$ of them, and adding velocities makes the
        grid impossible. Sampling-based planners store only the samples they
        draw, and the same code runs in any dimension.
  \item \rrt repeats the four primitives of \cref{def:ch16-primitives}:
        \textsc{Sample} (uniform in
        $\mathcal{X}$, the goal with probability $p_{\mathrm{goal}}$),
        \textsc{Nearest} (Euclidean nearest vertex), \textsc{Steer} (at most
        $\eta$ toward the sample) and \textsc{CollisionFree} (the whole
        segment, by subdivision with the margin $r+\delta/2$). A vertex that
        enters the goal region ends the search.
  \item The Voronoi bias explains the fast exploration: a vertex is
        \emph{selected} with probability equal to the relative volume of its
        Voronoi region (\cref{eq:ch16-voronoi}), and extended when the step is
        collision-free, so the tree is pulled into the largest unexplored
        regions.
  \item \rrt is probabilistically complete (\cref{thm:ch16-complete}) for any
        $\eta>0$ and $p_{\mathrm{goal}}<1$, with an expected number of
        iterations that grows like $(1/w)^{d}$ for a passage of width $w$.
        It is not optimal, and its best path converges almost surely to a
        suboptimal one (\cref{thm:ch16-suboptimal}); \rrtstar fixes this.
  \item \rrt-Connect grows two trees and connects them greedily; on the
        worked example it needs $2.4$ times fewer iterations than \rrt.
        Kinodynamic \rrt replaces \textsc{Steer} by forward simulation of
        sampled controls and needs a state-space metric for \textsc{Nearest}.
  \item Shortcutting and pruning never increase the length and remove most
        of the excess ($23.85\to18.39$ on the worked example) but cannot
        carry a path across an obstacle.
  \item Nearest-neighbor search is $\bigO{n}$ by a scan and $\bigO{\log n}$
        with a kd-tree; collision checks dominate in clutter. Fix the seed,
        report statistics over seeds, and never test only the vertices.
\end{itemize}
\end{summary}

\paragraph{Further reading.}
\rrt was introduced by LaValle in a technical report~\cite{lavalle1998rrt},
and the journal version with Kuffner~\cite{lavalle2001randomized} contains
the kinodynamic formulation and the completeness argument, which
\textcite{kleinbort2019completeness} revisited with a rigorous proof.
\rrt-Connect is due to Kuffner and LaValle~\cite{kuffner2000rrtconnect}.
LaValle's book~\cite{lavalle2006planning}, freely available online, treats
sampling-based planning in depth, including collision detection, metric
spaces and kinodynamic planning, and Choset et
al.~\cite{choset2005principles} give a compact textbook chapter. The
probabilistic roadmap of Kavraki et al.~\cite{kavraki1996prm} is the
multi-query counterpart. Karaman and Frazzoli~\cite{karaman2011sampling}
proved \cref{thm:ch16-suboptimal} and introduced \rrtstar, the subject of
\cref{ch:ch17} together with \irrtstar~\cite{gammell2014informed}. The
kd-tree is due to Bentley~\cite{bentley1975kdtree}, and Geraerts and
Overmars~\cite{geraerts2007creating} compare shortcutting with other path
improvement methods.

% ---------------------------------------------------------------------
\section{Exercises}
\label{sec:ch16-exercises}

\begin{exercise}[\difficulty{1}]
\label{exr:ch16-gridsizes}
A survey drone must plan through a $200\times200\times60$~m volume with an
accuracy of $0.25$~m. How many cells does the grid have, how much memory
does one byte per cell take, and how many neighbors does a cell have in the
26-connected lattice? Repeat for a state that also contains the velocity
($\pm4$~m/s per axis at $0.25$~m/s). Compare with the memory of an \rrt with
$20\,000$ vertices in three dimensions (three 8-byte coordinates and one
4-byte parent index per vertex).
\end{exercise}

\begin{exercise}[\difficulty{1}]
\label{exr:ch16-trace}
Continue \cref{tab:ch16-trace} by hand for iterations~11 and~12 of
\cref{ex:ch16-map}, given the samples $\state_{\mathrm{rand}}=(9.70,5.16)$ and
$(6.23,7.77)$: find the nearest vertex among vertices~0 to~6 (their
coordinates are in the table), compute $\state_{\mathrm{new}}$ from
\cref{eq:ch16-steer} with $\eta=0.5$, and decide whether the segment enters
the wall $[3,5]\times[0,6]$. Then explain, using the Voronoi bias, why
almost every sample in iterations~7 to~16 selects vertex~5 or~6.
\end{exercise}

\begin{exercise}[\difficulty{2}]
\label{exr:ch16-resolution}
(a)~Give a segment and a wall of thickness $0.04$ such that test points spaced
$\delta=0.05$ apart, tested \emph{without} the inflation by $\delta/2$, all
lie outside the wall although the segment crosses it. (b)~The margin
$r+\delta/2$ of \cref{thm:ch16-resolution} is sufficient but not the
smallest: show that the margin $\sqrt{r^{2}+\delta^{2}/4}$ also guarantees
a collision-free segment, and that no smaller margin does.
(c)~Let $d_0=\dist(\vect{a},\mathcal{O})-r$ be the free distance of the first
endpoint. Explain why no point of the segment within $d_0$ of $\vect{a}$ needs to
be tested, and design an adaptive check that uses this idea at every test
point (exact, no $\delta$, fast in open space).
\end{exercise}

\begin{exercise}[\difficulty{2}]
\label{exr:ch16-completeness}
This exercise completes \cref{thm:ch16-complete}. (a)~Let $T$ be the number
of iterations until the tree reaches the last ball $B_m$. Using the
binomial distribution of the number of samples that land in ``the next
ball'' among $n$ iterations, show that
$\Prob(T>n)\le\Prob\bigl(\mathrm{Bin}(n,p)<m\bigr)$ and deduce, with a
Chernoff bound, that the failure probability decays exponentially in $n$
once $n>2m/p$. (b)~Give a map on which \cref{alg:ch16-rrt} with
$p_{\mathrm{goal}}=1$ never finds a path although one exists, and explain
which step of the proof fails. (c)~Where exactly does the proof use
$\eta>0$, and what would go wrong for an algorithm whose step size shrinks
like $\eta_i=1/i$ at iteration $i$?
\end{exercise}

\begin{exercise}[\difficulty{2}]
\label{exr:ch16-voronoi}
Consider the tree with vertices $\vect{v}_0=(0,0)$, $\vect{v}_1=(1,0)$ and $\vect{v}_2=(2,0)$ in
the square $\mathcal{X}=[-1,3]\times[-2,2]$, with $p_{\mathrm{goal}}=0$.
Draw the Voronoi regions of the three vertices and compute the probability
that each one is chosen as $\state_{\mathrm{near}}$ by the next sample, using
\cref{eq:ch16-voronoi}. Then show in general that for a tree that is a single
straight chain of $k+1$ vertices spaced $\eta$ apart, centered in a large
square map of side $L$ with $L\gg k\eta$, each of the two tips is chosen with
probability close to $1/2$ and every interior vertex with probability at most
$\eta/L$. Finally explain why the picture changes when the chain starts near
a corner and points into the map, as in \cref{ex:ch16-map}: the region of the
single tip is then almost the whole square, so that tip takes almost all of
the probability.
\end{exercise}

\begin{exercise}[\difficulty{2} Coding]
\label{exr:ch16-goalbias}
Using \code{ch16\_rrt.py}, run \rrt on \cref{ex:ch16-map} for
$p_{\mathrm{goal}}\in\set{0,0.02,0.05,0.1,0.2,0.5,0.9}$ with $50$ seeds each
and plot the median and the quartiles of the number of iterations to the
first solution. Explain the shape of the curve with the arguments of
\cref{sec:ch16-pitfalls}. Then move the goal to $(9,1)$, where the straight
line from the start is blocked by the first wall only, and repeat: does the
best bias change?
\end{exercise}

\begin{exercise}[\difficulty{3}]
\label{exr:ch16-shortcut}
(a)~Prove that line~\ref{alg:ch16-shortcut:replace} of
\cref{alg:ch16-shortcut} never increases the length of the path and that
the result is collision-free whenever the input is. (b)~Prove that
\textsc{Prune} returns a path with at most as many vertices as the input and
give an example where its greedy jump of line~\ref{alg:ch16-shortcut:jump}
returns a longer path than a different choice of jumps would. (c)~Take the
wall $[4,6]\times[2,9.5]$ in $[0,10]^{2}$, the start $(1,5)$, the goal
$(9,5)$ and the path $(1,5),(1,9.75),(9,9.75),(9,5)$ over the top of the
wall. Prove that no sequence of shortcuts can produce a path that passes
below the wall, although that route is shorter.
\end{exercise}

\begin{exercise}[\difficulty{3} Coding]
\label{exr:ch16-coding}
This is the first half of the Week~8 coding exercise of the study plan
(\cref{ch:appA}); \cref{ch:ch17} completes it. Build a 3D obstacle field of about $20$ boxes
in a $50\times50\times20$~m volume (or reuse \code{scene\_3d()}) and plan
between random start--goal pairs with (a)~\rrt and \rrt-Connect from
\code{ch16\_rrt.py} followed by shortcutting, (b)~\rrtstar from
\cref{ch:ch17} or a library, and (c)~grid \astar from \cref{ch:ch04} on a
26-connected lattice at resolutions $2$, $1$ and $0.5$~m with inflated
obstacles. Over $30$ queries record success rate, path length relative to
the best found by any method, computation time and memory (vertices or
cells). Plot length against time and state in one paragraph, with your
numbers, when you would use graph search and when sampling: the Week~8
milestone of \cref{ch:appA}. Finally scale the map to
$200\times200\times60$~m at $0.25$~m and report which methods still run.
\end{exercise}
